% Transcription — fonds Alexandre Grothendieck, shelfmark 127, batch 2 (pages 21-23).
% Established from the facsimile archives/batches/127/batch-02.pdf, prepared
% from a copy of 127.pdf fetched through the project's relay (PDF page k =
% archivists' page 20+k, no cover sheet), per the skill
% transcribe-grothendieck. The folder is twenty-three pages; this is the
% second and last batch, three pages. Batch 1 (pages 1-20) is transcribed
% separately and in parallel by another pass and may not exist yet, so the
% notation below was fixed from these pages alone (they need nothing of
% folders 126 and 128 beyond the fibre-product and exactness conventions:
% no Amitsur complex or H^1 appears here).
%
% Pass: Opus 5.5 (claude-opus-5-5), 2026-09-26 — first pass, unchecked against the pages by a human.
%
% Dating: the folder is « s.d. » in src/content/catalogue.ts; its inventory
% group « Descentes » (cotes 126 to 133) is dated [avant 1970], and \dating
% says both, per the skill's group rule, exactly as folders 126 and 128 do.
% Nothing on these leaves dates them.
%
% Pages skipped: none. Page 21 is a tan sheet, otherwise blank, inscribed
% top right in brown ink « foncteurs sur / [struck: Morphismes + one word] /
% ramifiés, / étales ... » — a cover in his hand for the leaves that
% follow. By the rule of hand.md §5 it takes no \page marker; his words
% become the section heading, with a \note. (The margin of p. 22, in the
% same brown ink, reads « Avec morphismes non ramifiés », which supports
% « ramifiés » and suggests the struck word is « non ».)
%
% Pages summarised: none. No Murre or Levelt notes appear in this batch.
%
% Pages transcribed: 22-23, his, fast drafting in grey-black ink with
% later additions in brown ink (the diagonal marginal on p. 22, the
% inserted hypothesis « F(X) -> F(X_red) injectif », the marginal square
% and the underlinings on p. 23). His pagination: none visible.
%
% What the batch is about. A descent lemma for functors on finite
% preschemes: S locally noetherian, C the category of finite preschemes
% over S, F a contravariant set-valued functor on C with F(X) -> F(X_red)
% injective and the sheaf condition F(Y) -> F(X) => F(X x_Y X) for every
% EFFECTIVE epimorphism X -> Y; then the sequence is exact for EVERY
% epimorphism, proved by factoring an epimorphism into strict (effective)
% ones and inducting on the number of factors (p. 22, with a 2x3 diagram of
% the two steps). Corollary: for u surjective, F(Y) -> F(X) is injective
% (reduce to X, Y reduced, where u is an epimorphism). Special case
% (p. 23): A an abelian (? reading doubtful) preschema over S, F(S') the
% set of abelian sub-preschemas of A' = A x_S S'; for S' -> S an
% epimorphism, 0 -> F(S) -> F(S') => F(S'') exact, condition (i) coming
% from rigidity and (ii) from descent theory. A struck « Proposition »,
% renamed « Corollaire » with a struck parenthetical name, announces an
% analogous statement and breaks off.
%
% Notation fixed once for the batch: \mathcal{C} for his round C;
% \mathrm{Ob}\,\mathcal{C}; X_{\mathrm{red}}; X \times_Y X for his fibre
% products; \rightrightarrows for his double arrow. Words he underlines are
% \emph.
%
% Readings to check first: p. 21 the struck line; p. 22 the struck
% condition (i), the phrase « on peut prouver déjà », the last two lines
% (the second is interlinear and may run on to a lost continuation);
% p. 23 « abélien(s) » in the Proposition, the label on the right-hand
% vertical of the marginal square, and the whole of the struck ending.
% Apparatus: 14 \ill{}, 15 \uncertain{}. The double vertical arrows of
% the p. 22 diagram are drawn single (the renderer has no shift), with a \note.

\documentclass[11pt,a4paper]{article}
% Le préambule commun à toutes les transcriptions du fonds Grothendieck.
%
% Il tient en une idée : **l'incertitude se compose, elle ne se résout pas.**
% Une transcription qui lisse un mot douteux a détruit la seule chose pour
% laquelle on la faisait — un lecteur doit pouvoir savoir, à chaque endroit,
% ce qui était sur le feuillet et ce qui a été ajouté. Les six macros qui
% suivent sont donc l'appareil critique, et non de la décoration.
%
% Les mêmes commandes sont reconnues par `scripts/render.mjs`, qui produit la
% vue de lecture du site. Une macro ajoutée ici doit l'être aussi là-bas,
% faute de quoi le rendu échouera bruyamment — ce qui est voulu : un rendu
% silencieusement faux est pire qu'un rendu absent.

\ProvidesPackage{grothendieck}[2026/08/08 Transcriptions du fonds Grothendieck]

\usepackage{amsmath,amssymb,amsthm}
\usepackage{mathrsfs}
\usepackage{stmaryrd}
% Les diagrammes commutatifs sont partout dans ce fonds : c'est la langue
% même de la géométrie algébrique de Grothendieck, et une transcription qui
% les rendrait en prose ne transcrirait rien.
\usepackage{tikz-cd}
% Français par défaut — c'est la langue des feuillets — mais l'anglais est
% chargé aussi : la traduction et le résumé sont des documents anglais, et la
% typographie française (espaces avant les deux-points) y serait une faute.
% Ces documents basculent par \selectlanguage{english} en tête de corps.
\usepackage[english,french]{babel}
\usepackage{fontspec}
\usepackage{geometry}
\geometry{a4paper,margin=25mm,marginparwidth=28mm}
\usepackage{marginnote}
\usepackage{xcolor}
% ulem, not soul. `soul`'s \st and \ul are built on a character-by-character
% scan that XeLaTeX's font machinery breaks: the document fails with an
% undefined control sequence rather than degrading. ulem does the same two jobs
% and is Unicode-engine safe. `normalem` stops it hijacking \emph, which the
% transcriptions use for Grothendieck's own emphasis.
\usepackage[normalem]{ulem}
\usepackage{hyperref}
\hypersetup{colorlinks=true,linkcolor=black,urlcolor=black,pdfborder={0 0 0}}

% --- Métadonnées du lot -----------------------------------------------------
% Reprises telles quelles par le script de rendu, qui les affiche en tête de
% la vue de lecture. Elles fixent ce que le fichier prétend couvrir : sans
% elles, une transcription flotte sans point d'attache dans un fonds de seize
% mille feuillets.

\newcommand{\folder}[1]{\gdef\grothFolder{#1}}
\newcommand{\batch}[1]{\gdef\grothBatch{#1}}
\newcommand{\leaves}[2]{\gdef\grothFirst{#1}\gdef\grothLast{#2}}
\newcommand{\foldertitle}[1]{\gdef\grothFoldertitle{#1}}
\gdef\grothFolder{?}\gdef\grothBatch{?}\gdef\grothFirst{?}\gdef\grothLast{?}\gdef\grothFoldertitle{}

% --- Le résumé --------------------------------------------------------------
%
% Il ouvre la lecture modernisée, sous le titre « Résumé », et il est composé
% en petit. Ce n'est pas un ornement : c'est le seul passage du document qui
% ne soit pas des mathématiques, et un lecteur doit pouvoir le reconnaître
% avant de l'avoir lu — puis le sauter s'il connaît déjà le sujet, ou s'y
% arrêter s'il ne le connaît pas. Le filet qui le suit marque la reprise du
% corps.

\newenvironment{resume}
  {\small\setlength{\parskip}{0.45em}}
  {\par\medskip\hrule\medskip}

% --- Mots-clés ---------------------------------------------------------------
%
% En anglais, en clôture du résumé de la lecture modernisée : le vocabulaire
% moderne sous lequel chercher ce que les feuillets construisent. Le manifeste
% du site (`npm run manifest`) les extrait pour étiqueter la cote entière —
% cette ligne est la source unique des tags, il n'y a pas de fichier à côté.
\newcommand{\keywords}[1]{\par\smallskip\noindent
  {\footnotesize\textsc{Keywords} --- \textit{#1}}\par}

% --- L'appareil critique ----------------------------------------------------

% Le numéro de feuillet, en marge. C'est l'ancre qui permet de régler un
% désaccord sur une formule en regardant *un* feuillet plutôt qu'une chemise
% de six cents. Le site s'en sert aussi pour tourner les pages du fac-similé
% quand on fait défiler la transcription.
\newcommand{\leaf}[1]{%
  \marginnote{\footnotesize\color{gray!70}\textsf{#1}}%
  \label{leaf:#1}\ignorespaces}

% Illisible : on ne devine pas. Un mot inventé qui se lit comme les autres est
% le pire résultat possible.
\newcommand{\ill}{\textcolor{red!60!black}{[\,\ldots\,]}}

% Lecture douteuse : le mot est donné, et signalé comme incertain.
\newcommand{\uncertain}[1]{\uline{#1}}

% Ajout de l'éditeur — une lettre manquante, un mot sous-entendu.
\newcommand{\add}[1]{\textcolor{blue!50!black}{[#1]}}

% Biffé par Grothendieck lui-même. Ce qu'il a rayé fait partie du document :
% une rature dit souvent où la pensée a changé de direction.
\newcommand{\struck}[1]{\sout{#1}}

% Note du transcripteur, sur l'état matériel du feuillet.
\newcommand{\note}[1]{\par\smallskip{\small\color{gray!60!black}\textit{[#1]}}\par\smallskip}

% Note marginale de Grothendieck, distincte du corps du texte.
\newcommand{\marginal}[1]{\par\smallskip\noindent\hspace{1em}%
  {\small\color{gray!60!black}#1}\par\smallskip}

% --- Titre ------------------------------------------------------------------

\AtBeginDocument{%
  \begin{center}
    {\large\bfseries Fonds Alexandre Grothendieck --- cote n\textsuperscript{o}~\grothFolder}\\[2pt]
    {\itshape\grothFoldertitle}\\[4pt]
    {\small feuillets \grothFirst--\grothLast\ (lot \grothBatch)}\\[2pt]
    {\footnotesize\color{gray!60!black}Transcription établie d'après le fac-similé numérisé
     par l'Université de Montpellier.\\
     Les lectures incertaines sont soulignées, les illisibles notées [\,\ldots\,],
     les ajouts entre crochets.}
  \end{center}
  \vspace{1em}}

\endinput


\folder{127}
\batch{2}
\pages{21}{23}
\dating{s.d. — le groupe « Descentes » (126 à 133) est daté [avant 1970]}
\watermark{Édition de démonstration}
\foldertitle{Descente non plate : notes manuscrites (s.d.)}

\begin{document}

\section*{Foncteurs sur \struck{Morphismes \uncertain{non}} \uncertain{ramifiés}, étales \ldots}
\note{ces mots sont les siens, à l'encre brune, en haut à droite de la
page 21, feuillet par ailleurs blanc qui sert de chemise aux pages
suivantes ; le mot biffé après « Morphismes » est peu lisible, et la note
en marge de la page 22 (« Avec morphismes non ramifiés ») appuie les
lectures « non » et « ramifiés »}

\page{22}

\marginal{Avec morphismes non ramifiés}
\note{en diagonale dans le coin supérieur gauche, à l'encre brune, séparé
du texte par un trait}

\emph{Lemme.} Soit $S$ un préschéma \struck{\ill{}} noeth., $\mathcal{C}$ la
catégorie des préschémas finis sur $S$, et $F$ un foncteur contravariant de
$\mathcal{C}$ dans la catégorie des ensembles. On suppose
\struck{(i) Si \ill{} $X \to Y$ induit un isomorphisme}
Pour tout $X \in \mathrm{Ob}\,\mathcal{C}$,
$F(X) \to F(X_{\mathrm{red}})$ est injectif
\note{la condition (i) primitive est biffée ; « Pour tout
$X \in \mathrm{Ob}\,\mathcal{C}$, $F(X) \to F(X_{\mathrm{red}})$ est
injectif » est ajouté en interligne au-dessus, « injectif » à l'encre
brune}

(ii) Pour tout épimorphisme \emph{effectif} $X \to Y$, \uncertain{on a}
une suite \emph{exacte}
\[
  F(Y) \to F(X) \rightrightarrows F(X \times_Y X) .
\]
\note{devant (ii) un signe noirci ; « effectif » n'est pas souligné, mais
écrit d'une autre encre}

Alors pour \emph{tout} épi\emph{morphisme} $X \to Y$, on a une suite exacte
\[
  F(Y) \to F(X) \rightrightarrows F(X \times_Y X) .
\]

\emph{Démonstration.} On utilise le fait que
$u : X \to Y$ peut
se \uncertain{factoriser} en produit d'épimorphismes stricts, \uncertain{on peut}
\uncertain{prouver} déjà que $F(Y) \to F(X)$ est injectif (grâce à (ii)). Pour
l'exactitude totale, récurrence sur le nombre de facteurs nécessaires ; si
$n = 1$, c'est (ii), supposons $n > 1$ et prouvé pour $n' < n$. On a donc
\[
  X \xrightarrow{\ u'\ } X'_1 \xrightarrow{\ u''\ } Y ,
\]
où $u'$ est produit de $n - 1$ épim.\ effectifs, et $u''$ est un épim.\
effectif. \uncertain{Mais} \ill{} partons de $\xi \in F(X)$ tel que ses deux
images dans $F(X \times_Y X)$ soient égales, les deux images dans
$F(X \times_{X'} X)$ sont \uncertain{a fortiori} égales, donc par hyp.\ de
récurrence $\xi$ provient d'un $\xi' \in F(X')$, reste à prouver que
\ill{} \ill{} \ill{} d'un $\xi \in F(X)$, on \ill{} que ses deux images
dans $F(X' \times_Y X')$ sont les mêmes. Cela résulte du fait que
$F(X' \times_Y X') \to F(X \times_Y X)$ est \emph{injectif},
\uncertain{lui-même} \uncertain{conséquence} \ill{} \uncertain{ceci}
\note{le texte écrit $X'_1$ dans la factorisation, puis $X'$ dans la
suite et dans le diagramme ; la dernière ligne, serrée sous la
précédente au bas du feuillet, s'arrête là}

\begin{tikzcd}[column sep=small, row sep=small, nodes={font=\scriptsize}]
F(Y) \arrow[r] & F(X') \arrow[r] \arrow[d] & F(X) \arrow[d] & \\
 & F(X' \times_Y X') \arrow[r] & F(X \times_Y X) \arrow[r] & F(X \times_{X'} X)
\end{tikzcd}

\note{diagramme en marge gauche, à hauteur de l'argument ; $\xi'$ est écrit
au-dessus de $F(X')$ et $\xi$ au-dessus de $F(X)$ ; les flèches verticales
sont doubles sur la page (deux flèches parallèles), simplifiées ici}

\page{23}

\emph{Corollaire.} Les conditions (i) (ii) impliquent que si
$u : X \to Y$ est \emph{surjectif}, alors $F(Y) \to F(X)$ est
\emph{injectif}. En effet, grâce à (i) on peut supposer $X$, $Y$ réduits,
mais alors $u$ est un épim., et nous avons déjà vu que \struck{\ill{}}
alors $F(Y) \to F(X)$ est injectif. OK.
\note{« Corollaire » et les soulignements de « surjectif » et « injectif »
sont à l'encre brune}

\begin{tikzcd}[column sep=small, row sep=small, nodes={font=\scriptsize}]
F(X) \arrow[d] & F(Y) \arrow[l] \arrow[d] \\
F(X_{\mathrm{red}}) & F(Y_{\mathrm{red}}) \arrow[l]
\end{tikzcd}

\note{carré en marge gauche, à l'encre brune ; la flèche verticale de
droite porte une marque, \uncertain{inj}, peu lisible}

Cas particulier

\emph{Proposition.} Soient $A$ préschéma \uncertain{abélien} sur $S$ loc.\
noeth., et pour tout $S'$ sur $S$, soit $F(S')$ l'ens.\ des
sous-préschémas \uncertain{abéliens} de $A' = A \times_S S'$. Alors si $S'$
est tel que $S' \to S$ soit un épimorphisme, on a une suite exacte
(où $S'' = S' \times_S S'$)
\[
  0 \to F(S) \to F(S') \rightrightarrows F(S'')
\]

En effet, la condition (i) du Lemme résulte de la rigidité, la condition
(ii) résulte de la théorie de la descente.

\struck{Proposition} Corollaire. (\struck{\ill{}} \ill{})
\struck{Supposons que $S$ soit} Énoncé analogue pour les \ill{}
\struck{le spectre d'}un \ill{} abéliens.
\note{« Proposition », souligné, est biffé et remplacé au-dessus par
« Corollaire » ; dans la parenthèse un nom biffé et un autre écrit
au-dessus, illisibles tous deux ; « Énoncé analogue pour les \ldots » est
écrit au-dessus de la ligne biffée ; l'énoncé s'arrête là}

\end{document}
